\section[Weighted tree ownership]{A charged weighted ownership decomposition}
\label{sec:op3-weighted-tree-ownership}

\begin{theorem}[Owned-load tree decomposition; \statusproved]
\label{thm:op3-weighted-tree-ownership}
Let $T$ be a supplied tree on $n\geq1$ vertices, with nonnegative
vertex loads $a_v$ and a threshold $\theta>0$. Write
$A_{\rm tot}=\sum_v a_v$. In $O(n)$ exact-real word work and total
allocated words, one can construct connected pieces $W_1,\ldots,W_h$
and an owner map $o:V\to\{1,\ldots,h\}$ such that:
\begin{enumerate}[label=(\roman*),leftmargin=*]
\item the pieces cover $V$, their tree edges partition $E(T)$, and any
two distinct pieces intersect in at most one vertex;
\item $v\in W_{o(v)}$, and each vertex has exactly one owner;
\item each piece has at most two vertices occurring in other pieces;
\item every nontrivial piece has owned load
$\sum_{v:o(v)=i}a_v<2\theta$; singleton owner pieces may be heavier;
\item either $h=1$ or $h\leq8A_{\rm tot}/\theta$.
\end{enumerate}
The construction makes no lower-bound assumption on a positive load or
on its ratio to another load. Every scan, list splice, refinement,
membership, allocation, temporary copy and output word is charged.
\end{theorem}
\begin{proof}
Root $T$ arbitrarily and process it in postorder. An unfinished child
bag is a connected subtree containing the child, together with the
vertices whose load it owns. Its owned mass is strictly below $\theta$.
Represent its edges and owners by two linked lists.

At a vertex $v$ with $a_v<\theta$, begin a pending bag owning $v$.
If $a_v\geq\theta$, first emit the singleton owner piece $\{v\}$
and begin an empty-mass bag with connector $v$ instead. Attach each
child bag by its parent edge and splice its two lists into the pending
bag. As soon as its mass reaches $\theta$, emit that connected piece
and reset the pending bag to connector $v$ with no second ownership.
Before a merge both masses were below $\theta$, so every nontrivial
flushed piece has mass in $[\theta,2\theta)$. Return the remaining
bag to the parent. At the root, emit the final remainder if it has an
edge or an owner; omit a redundant empty connector.

Every vertex is owned once and every tree edge enters exactly one
emitted piece. The only repeated vertex after a flush is its connector.
The pieces are connected and edge-disjoint, so an intersection containing
two vertices would duplicate their unique tree path, which is impossible.
All emitted pieces except possibly the root remainder own mass at least
$\theta$. Thus the initial number $h_0$ satisfies
$h_0\leq A_{\rm tot}/\theta+1$; nontrivial initial pieces have owned
mass below $2\theta$. A heavy singleton has no tree edge.

It remains to limit boundary count without losing ownership. Let $q$
be the number of shared vertices of the initial decomposition and $I$
their total number of piece incidences. Each tree piece has one more
vertex than edge, including singletons, so
\[
 \sum_i|W_i|=n-1+h_0,\qquad I=h_0+q-1.
\]
If $h_0>1$, every piece meets another one. Since every shared vertex
has at least two incidences, $q\leq h_0-1$ and $I\leq2(h_0-1)$.

Keep a piece with at most two boundaries intact. For a piece with
$b>2$ boundaries, peel its nonboundary leaves to obtain the minimal
Steiner tree spanning those boundaries. Mark the boundaries and the
Steiner branching vertices. There are at most $b-2$ branching vertices
outside the boundary set, hence at most $t\leq2b-2$ marked vertices.
Make a piece from each maximal corridor between marked vertices,
including every hanging subtree attached to an interior corridor vertex.
At each marked vertex, group all its hanging branches into one additional
piece if any exist. There are at most $(t-1)+t\leq4b-5$ new pieces.
Every corridor piece has only its two marked endpoints as possible
boundaries; every additional hanging piece has only its marked root.

These pieces partition the old piece's tree edges and meet only at
marked vertices. Assign every old owned vertex to one new piece
containing it. No new piece acquires load from another old piece, so
the strict $2\theta$ bound is preserved. Heavy singleton pieces are
kept intact. When $h_0>1$, every old piece has at least one boundary,
and its number of replacements is at most four times that boundary
count. Therefore
\[
 h\leq4I\leq8(h_0-1)\leq8A_{\rm tot}/\theta.
\]
If $h_0=1$, no refinement is necessary and $h=1$.

For the work bound, linked edge/owner lists are spliced in constant
work; a subtree is never copied up its ancestor chain. There is one
list node per tree edge and per vertex owner, and only $O(n)$ bag
headers: a nontrivial flush owns a previously unassigned tree edge,
while there is at most one heavy singleton per vertex. Materializing
initial membership lists reads each edge and owner once, using one
reusable stamp array. Their total size is $n-1+h_0=O(n)$.

Build each old piece's local adjacency from its own edges. A reusable
global-to-local index array is touched only at its members. Peeling,
reverse attachment recovery, corridor traversal and hanging-edge
assignment each scan those local records a constant number of times.
Each old tree edge is copied into exactly one refined piece; each owner
is reassigned once. The final membership size is $n-1+h=O(n)$, because
every nontrivial final piece contains an edge and there are at most $n$
singleton owner pieces. Reusable global arrays and all local temporary
arrays, including their release work, total $O(n)$ over the call.
Retaining both initial and final piece lists still occupies $O(n)$
words; no discarded history is omitted from the allocation count.
\end{proof}

\begin{corollary}[Assigning original edge loads; \statusproved]
\label{cor:op3-weighted-edge-ownership}
Let an additional supplied graph have positive edge loads $w_e$ on the
same vertices, with $W=\sum_e w_e>0$. Set
$a_v=\sum_{e\ni v}w_e$ and apply
\cref{thm:op3-weighted-tree-ownership} with
$\theta=64W/j$, where $64\leq j\leq n$.
Assign each original edge to its one or two endpoint owner pieces.
The refined decomposition has at most $j/4$ pieces, and every
nontrivial piece is assigned original edge load below $128W/j$.
The additional load assembly and endpoint assignment cost $O(m+n)$.
\end{corollary}
\begin{proof}
Total vertex load is $2W$, so the piece-count bound is
$\max\{1,16W/\theta\}=j/4$. An original edge assigned to a piece
has at least one endpoint owned there, so its load is included at least
once in that piece's owned vertex load. Thus the assigned edge load is
at most its owned load, which is below $2\theta=128W/j$ for nontrivial
pieces. Every endpoint and ownership record is scanned once.
\end{proof}

\paragraph{Audit and composition.}
The implementation \texttt{weighted\_tree\_ownership.py} passes 9,042
exact supplied-tree cases, including all unlabeled trees through nine
vertices with every root and three load profiles, threshold ties,
zero loads, heavy singletons, and larger stars, paths and random trees
through 512 vertices. It validates 127,547 pieces and 1,483,000 pairwise
intersections, and records 10,486,152 charged construction units.
There are 6,562 actual Steiner refinements. The large-ratio profile
contains $2^{-80}$ and $2^{80}$ without rounding. Validation graphs and
pairwise intersection scans are independent audit work, not construction
cost. Hashes and exact counts are in
\texttt{WEIGHTED\_TREE\_OWNERSHIP\_AUDIT.json}.

This supplies the decomposition and ownership part of the weighted
constructor. The charged tree-distance/congestion computation and the
additional argument converting owned load into routed stretch are proved
in \cref{sec:op3-weighted-corridor-routing}. Source/confidence composition
still remains; the supplied near-linear recurrence is \statusconditional.
